% Schemi pronti da copiare negli src/. Questo file NON viene compilato.
% Figure e tabelle usano [H]: restano nel punto del testo in cui sono scritte.

% ——— Tabella di misure ———
\begin{table}[H]
    \centering
    \begin{tabular}{ccc}
        \toprule
        $x$ (u.m.) & $y$ (u.m.) & $\sigma_y$ (u.m.) \\
        \midrule
        1.0 & 2.0 & 0.1 \\
        \bottomrule
    \end{tabular}
    \caption{CAPTION}
    \label{tab:misure1}
\end{table}

% ——— Tabella di confronto con Z-test (numeri da results.tex) ———
\begin{table}[H]
    \centering
    \begin{tabular}{lcccc}
        \toprule
        Grandezza & Misurato & Atteso & $Z$ & Esito \\
        \midrule
        $b$ (u.m.) & \res{p1b} & \res{p1b:ref} & \res{p1b:z} & compatibile \\
        \bottomrule
    \end{tabular}
    \caption{CAPTION}
    \label{tab:risultati1}
\end{table}

% ——— Risultato nel testo ———
Si ottiene $b = (\res{p1b})$~u.m.

% ——— Risultati finali in riga (conclusioni) ———
\[
\begin{aligned}
  &b = (\res{p1b})\ \text{u.m.} && Z = \res{p1b:z}
\end{aligned}
\]

% ——— Grafico singolo ———
\begin{figure}[H]
    \centering
    \includegraphics[width=0.9\linewidth]{plots/part1.png}
    \caption{CAPTION}
    \label{plt:part1}
\end{figure}

% ——— Due grafici o foto affiancati ———
\begin{figure}[H]
    \centering
    \begin{subfigure}[b]{0.48\textwidth}
        \includegraphics[width=\linewidth]{plots/part1.png}
        \caption{CAPTION}
        \label{plt:a}
    \end{subfigure}\hfill
    \begin{subfigure}[b]{0.48\textwidth}
        \includegraphics[width=\linewidth]{plots/part2.png}
        \caption{CAPTION}
        \label{plt:b}
    \end{subfigure}
    \caption{CAPTION}
\end{figure}

% ——— Pagina di soli grafici a margini ridotti ———
\clearpage
\thispagestyle{empty}
\newgeometry{top=1cm, bottom=0.5cm, left=0.5cm, right=0.5cm}
% ... figure ...
\restoregeometry
\clearpage

% ——— Schema di un circuito (circuitikz) ———
% Richiede \usepackage[european]{circuitikz} in styles.tex.
\begin{figure}[H]
    \centering
    \begin{circuitikz}
        \draw (0,0) to[sV, l=$v_i$] (0,2)        % generatore di continua: to[battery1, l=$v_i$]
              to[C, l=$C$] (3,2)
              to[R, l=$R$, *-*] (3,0) -- (0,0);
        \draw (3,2) to[short, -o] (4.5,2) node[right] {$v_o$};
        \draw (3,0) to[short, -o] (4.5,0);
        \draw (1.5,0) node[ground] {};
    \end{circuitikz}
    \caption{CAPTION}
    \label{fig:circuito1}
\end{figure}
